\documentclass[11pt,a4paper]{article}

\usepackage[margin=1in]{geometry}
\usepackage{amsmath,amssymb,amsthm}
\usepackage{mathtools}
\usepackage{hyperref}
\usepackage[capitalise,noabbrev]{cleveref}

\newtheorem{theorem}{Theorem}
\newtheorem{lemma}[theorem]{Lemma}
\newtheorem{corollary}[theorem]{Corollary}
\newtheorem{proposition}[theorem]{Proposition}
\theoremstyle{definition}
\newtheorem{definition}[theorem]{Definition}
\newtheorem{game}[theorem]{Game}
\theoremstyle{remark}
\newtheorem{remark}[theorem]{Remark}

\newcommand{\KEM}{\ensuremath{\mathsf{KEM}}}
\newcommand{\KGen}{\ensuremath{\mathsf{KGen}}}
\newcommand{\Encap}{\ensuremath{\mathsf{Encap}}}
\newcommand{\Decap}{\ensuremath{\mathsf{Decap}}}
\newcommand{\KDF}{\ensuremath{\mathsf{KDF}}}
\newcommand{\Adv}{\ensuremath{\mathbf{Adv}}}
\providecommand{\Exp}{}\renewcommand{\Exp}{\ensuremath{\mathbf{Exp}}}
\providecommand{\Game}{}\renewcommand{\Game}{\ensuremath{\mathbf{G}}}
\newcommand{\negl}{\ensuremath{\mathsf{negl}}}
\newcommand{\xwing}{\textsc{X-Wing}}
\newcommand{\zwing}{\textsc{Z-Wing}}
\newcommand{\zap}{\textsc{ZAP}}
\newcommand{\mlkem}{\ensuremath{\mathsf{ML\text{-}KEM}}}
\newcommand{\xtf}{\ensuremath{\mathsf{X25519}}}
\newcommand{\sha}{\ensuremath{\mathsf{SHA3}}}
\newcommand{\hkdf}{\ensuremath{\mathsf{HKDF}}}
\newcommand{\BB}{\ensuremath{\{0,1\}}}

\title{Security Proofs for \zwing{} in the \zap{} Handshake}
\author{\zap{} Protocol Authors\\\texttt{zap-proto.io}}
\date{\today}

\begin{document}
\maketitle

\begin{abstract}
This document accompanies the paper ``Post-Quantum Hybrid KEMs in the
\zap{} Transport Handshake'' and contains the full proofs of the IND-CCA2
security of \zwing{} and the IND-CCA2-AKE security of the \zap{} handshake
instantiated with \zwing{}. We work in the random-oracle model and use a
sequence-of-games style following \cite{shoup2004games}.
\end{abstract}

\tableofcontents

\section{Notation and Preliminaries}\label{sec:notation}

We let $\lambda$ be the security parameter and write $\negl(\lambda)$
for any function negligible in $\lambda$. ``QPT'' means quantum
polynomial-time; the random oracle is in the QROM where the analysis
requires it. For a finite set $S$, $x \xleftarrow{\$} S$ denotes uniform
sampling. For an algorithm $A$, $y \gets A(x)$ denotes execution.

\subsection{KEM syntax and IND-CCA2}

A KEM $\KEM = (\KGen, \Encap, \Decap)$ consists of:
\begin{itemize}
\item $\KGen(1^\lambda) \to (pk, sk)$,
\item $\Encap(pk) \to (c, ss)$ with $ss \in \BB^{\ell}$,
\item $\Decap(sk, c) \to ss$ or $\bot$.
\end{itemize}
\emph{Correctness}: $\Pr[\Decap(sk, c) = ss : (pk, sk) \gets \KGen,\ (c, ss) \gets \Encap(pk)] \ge 1 - \negl(\lambda)$.

\begin{definition}[IND-CCA2]\label{def:indcca}
For $\KEM$ and adversary $\mathcal{A}$,
$\Adv^{\mathsf{ind\text{-}cca}}_{\KEM}(\mathcal{A}) =
 |\Pr[\Exp^{\mathsf{ind\text{-}cca}, 1}_\KEM(\mathcal{A}) = 1]
 - \Pr[\Exp^{\mathsf{ind\text{-}cca}, 0}_\KEM(\mathcal{A}) = 1]|$,
where $\Exp^{\mathsf{ind\text{-}cca}, b}_\KEM$:
\begin{enumerate}
\item $(pk, sk) \gets \KGen(1^\lambda)$;
\item $(c^*, ss_0) \gets \Encap(pk)$; $ss_1 \xleftarrow{\$} \BB^\ell$;
\item give $\mathcal{A}$ access to $\mathcal{O}_{\Decap}(\cdot)$, refusing $c^*$;
\item $b' \gets \mathcal{A}(pk, c^*, ss_b)$; output $b'$.
\end{enumerate}
\end{definition}

\subsection{Hybrid KEM and freshness}

\begin{definition}[Hybrid freshness]\label{def:fresh}
For a session $\pi$ using \zwing{} with components $(\KEM_1, \KEM_2, \KEM_3)$,
$\pi$ is \emph{fresh} if there exists $i^\star \in \{1, 2, 3\}$ such that
neither party's $\KEM_{i^\star}$ secret in $\pi$ has been revealed via
$\mathsf{StaticReveal}$ or $\mathsf{EphemeralReveal}$, and the test
session's transport key has not been revealed.
\end{definition}

The freshness condition is the relaxation that distinguishes hybrid AKE from
single-primitive AKE: an adversary may compromise up to two of the three
KEM secrets and the protocol still must hide the session key.

\section{Lemma 1: KDF Combiner Robustness}\label{sec:lem1}

\begin{lemma}[Robust combiner]\label{lem:combiner}
Let $H : \BB^* \to \BB^{n}$ be modelled as a random oracle, with output
length $n = 256$. Define
\[
F(ss_1, ss_2, ss_3, c, pk) := H(\texttt{"Z-Wing-v1"} \parallel ss_1 \parallel ss_2 \parallel ss_3 \parallel c \parallel pk).
\]
For any adversary $\mathcal{A}$ making at most $q_H$ queries to $H$,
\begin{enumerate}
\item if for some $i^\star \in \{1, 2, 3\}$ the value $ss_{i^\star}$ is
uniformly random and unknown to $\mathcal{A}$, then $F(\cdot)$ is
$2^{-n}$-statistically close to uniform conditioned on
$\mathcal{A}$'s view of $(ss_j)_{j \neq i^\star}$ and $(c, pk)$;
\item the only way $\mathcal{A}$ can distinguish $F(\cdot)$ from uniform is
to query $H$ on the exact preimage including $ss_{i^\star}$, an event with
probability at most $q_H \cdot 2^{-|ss_{i^\star}|}$;
\item $F$ binds the ciphertext $c$ and public key $pk$: for any
$(ss_1, ss_2, ss_3) \neq (ss'_1, ss'_2, ss'_3)$ or
$(c, pk) \neq (c', pk')$, $\Pr[F(\ldots) = F'(\ldots)] \le 2^{-n}$.
\end{enumerate}
\end{lemma}

\begin{proof}
We prove each clause.

\textbf{(1) and (2): Hidden-input randomness.}
Without loss of generality, suppose $ss_{i^\star}$ is uniformly random in
$\BB^{\ell_{i^\star}}$ and unknown to $\mathcal{A}$. The output of $H$ at
input $X = \texttt{"Z-Wing-v1"} \parallel ss_1 \parallel ss_2 \parallel ss_3 \parallel c \parallel pk$
is, in the ROM, a uniformly random string conditioned on the event $E$ that
$\mathcal{A}$ has not queried $H$ at $X$. Since $X$ contains $ss_{i^\star}$
which has $\ge \lambda$ bits of min-entropy from $\mathcal{A}$'s view (we
assume $\ell_{i^\star} \ge 256$), the probability that $\mathcal{A}$ guesses
$X$ in any one of its $q_H$ oracle queries is at most $q_H / 2^{\ell_{i^\star}}$.
Conditioned on $\bar{E}$ (no such query), $H(X)$ is uniform in $\BB^n$ from
$\mathcal{A}$'s view. Hence
\[
\Delta\bigl(F(\cdot);\ U_n \mid \mathcal{A}\bigr) \le \frac{q_H}{2^{\ell_{i^\star}}} + 2^{-n} \le \frac{q_H + 1}{2^n}
\]
since $\ell_{i^\star} \ge n = 256$.

\textbf{(3): Binding.}
$F$ takes a unique input string for distinct $(ss_1, ss_2, ss_3, c, pk)$
because the components are length-tagged or fixed-length (the public-key
and ciphertext lengths of $\KEM_1, \KEM_2, \KEM_3$ are fixed by the algorithm
choice; the label is fixed; and the shared secrets are 32 bytes each).
The probability that the random oracle outputs the same value at two distinct
inputs is, by a union bound over $q_H$ queries, at most
$\binom{q_H}{2} 2^{-n}$.
\qedhere
\end{proof}

\begin{remark}
Clause (3) is what closes the ``ciphertext malleability'' gap of naive
combiners. If we omitted $c$ or $pk$ from $H$'s input, an adversary could
maul one component's ciphertext (e.g.\ a CCA-malleable but not CCA-secure
component) without changing the derived key, breaking IND-CCA2 of the hybrid.
\end{remark}

\section{Lemma 2: Transcript Binding Prevents Key Reuse}\label{sec:lem2}

\begin{lemma}[Transcript binding]\label{lem:transcript}
Let $\tau$ be the \zap{} transcript and let $K_{\mathsf{tx}} = \hkdf(ss, \tau)$
be the transport-key derivation. For two sessions $\pi, \pi'$ with derived
keys $K, K'$, $K = K'$ implies $\tau = \tau'$ and $ss = ss'$
except with probability $q_H \cdot 2^{-n}$ over the random oracle modelling
of $\hkdf$, where $q_H$ is the total number of oracle queries.
\end{lemma}

\begin{proof}
$\hkdf$-Extract followed by $\hkdf$-Expand is a two-stage extract-then-expand
construction. We model the inner HMAC-based extract and the inner expand as
independent random oracles $H_e, H_x$ (this is conservative; the actual
HKDF analysis \cite{krawczyk2010hkdf} uses a weaker assumption).

The transcript $\tau$ contains a fresh nonce per session
($|\mathsf{nonce}| = 256$ in \zap{}) and the public-key fingerprints of both
peers. For two sessions $\pi, \pi'$, the probability that
$\tau = \tau'$ is bounded by the collision probability of the nonces, namely
$2^{-256}$ over an honest randomness sampling (a single pair).

If $\tau \neq \tau'$, then $H_x(H_e(0, ss) \parallel \tau) = H_x(H_e(0, ss') \parallel \tau')$
forces a collision on $H_x$, occurring with probability at most
$\binom{q_H}{2} 2^{-n}$.

If $\tau = \tau'$ but $ss \neq ss'$, then $H_e(0, ss) \neq H_e(0, ss')$ except
with collision probability $\binom{q_H}{2} 2^{-n}$, after which the same
$H_x$ collision argument applies.

Combining, $\Pr[K = K' \wedge (\tau, ss) \neq (\tau', ss')] \le \binom{q_H}{2} 2^{-n}$.
\qedhere
\end{proof}

\section{Theorem 1: Z-Wing IND-CCA2}\label{sec:thm1}

\begin{theorem}[\zwing{} IND-CCA2]\label{thm:zw-cca}
For any QPT adversary $\mathcal{A}$ against the IND-CCA2 security of \zwing{},
making at most $q_H$ random-oracle queries and $q_D$ decapsulation queries,
\[
\Adv^{\mathsf{ind\text{-}cca}}_{\zwing}(\mathcal{A}) \le
\min_{i \in \{1,2,3\}} \Adv^{\mathsf{ind\text{-}cca}}_{\KEM_i}(\mathcal{B}_i) +
\frac{q_H + q_D + 1}{2^n} +
\delta_{\KEM_1} + \delta_{\KEM_2} + \delta_{\KEM_3},
\]
where $\delta_{\KEM_i}$ is the correctness error of $\KEM_i$ and the running
time of each $\mathcal{B}_i$ is approximately $t(\mathcal{A}) + (q_H + q_D) \cdot \mathsf{poly}(\lambda)$.
\end{theorem}

\begin{proof}
We prove the bound for the case $i^\star = 2$ (the $\mlkem$-768 component);
the cases $i^\star = 1$ and $i^\star = 3$ are symmetric.

\paragraph{Game $\Game_0$.} Real IND-CCA2 game against \zwing{}.
By definition,
$\Pr[\Game_0 \to 1] = \frac{1}{2} + \frac{1}{2} \Adv^{\mathsf{ind\text{-}cca}}_{\zwing}(\mathcal{A})$.

\paragraph{Game $\Game_1$.} Identical to $\Game_0$ except that the simulator
maintains a list $\mathcal{L}_H$ of all queries to $H$ and aborts if
$\mathcal{A}$ ever submits to its decapsulation oracle a ciphertext $c$
whose decapsulation produces $(ss_1, ss_2, ss_3)$ such that the input
$X = \texttt{label} \parallel ss_1 \parallel ss_2 \parallel ss_3 \parallel c \parallel pk$
has not been queried to $H$.

\emph{Bounding $|\Pr[\Game_0 \to 1] - \Pr[\Game_1 \to 1]|$.}
The abort happens only when the simulator must answer a decap query whose
$H$-input has not been pre-queried. In the random-oracle model, the value of
$H(X)$ for an unqueried $X$ is uniform from $\mathcal{A}$'s view. Therefore
the simulator can defer answering: it samples a random $ss$, returns it, and
``programs'' $H(X) = ss$ later. The only failure case is if $\mathcal{A}$
queries $H$ at $X$ before submitting the decap query; the probability of
$\mathcal{A}$ pre-querying $X$ for an honest decap output is bounded by
$q_H q_D / 2^n$ (each of the $q_D$ decap inputs picks a random
$ss_2$, hence a random $X$; collision with any of $q_H$ queries has
probability $q_H/2^n$). Hence
\[
|\Pr[\Game_0 \to 1] - \Pr[\Game_1 \to 1]| \le \frac{q_H q_D}{2^n}.
\]

\paragraph{Game $\Game_2$.} Identical to $\Game_1$ except in the challenge
ciphertext: the value $ss_2^*$ used by the simulator inside $H$ is
\emph{replaced by a uniformly random string} $\widetilde{ss_2} \xleftarrow{\$} \BB^{\ell_2}$.
The decap oracle remains honest with respect to $sk_2$, and any decap query
that decapsulates to a value other than $\widetilde{ss_2}$ is answered honestly.

\emph{Bounding $|\Pr[\Game_1 \to 1] - \Pr[\Game_2 \to 1]|$.}
Construct $\mathcal{B}_2$, an IND-CCA2 adversary against $\KEM_2 = \mlkem$-768:
$\mathcal{B}_2$ receives $(pk_2, c_2^*, ss_2^*)$ and must guess whether
$ss_2^*$ is real or random. $\mathcal{B}_2$ generates $(pk_1, sk_1), (pk_3, sk_3)$
itself, sets $pk = (pk_1, pk_2, pk_3)$, simulates the random oracle for
$\mathcal{A}$, and answers decap queries by using $sk_1, sk_3$ for the
classical and $\KEM_3$ components and forwarding the $\KEM_2$ component to
its own decap oracle. The challenge \zwing{} ciphertext is
$c^* = (pk_1', c_2^*, c_3^*)$ where $\mathcal{B}_2$ samples $pk_1', c_3^*$
honestly. The simulated key $ss^* = H(X)$ uses the $ss_2^*$ from
$\mathcal{B}_2$'s challenger.
If $ss_2^*$ is real, the simulation is exactly $\Game_1$; if random, it is
$\Game_2$. Hence
\[
|\Pr[\Game_1 \to 1] - \Pr[\Game_2 \to 1]| \le \Adv^{\mathsf{ind\text{-}cca}}_{\KEM_2}(\mathcal{B}_2).
\]

\paragraph{Game $\Game_3$.} Identical to $\Game_2$ except that $ss^*$, the
challenge \zwing{} shared secret, is replaced by a uniformly random
$\widetilde{ss^*} \xleftarrow{\$} \BB^n$.

\emph{Bounding $|\Pr[\Game_2 \to 1] - \Pr[\Game_3 \to 1]|$.}
By Lemma~\ref{lem:combiner}, with $i^\star = 2$ and $\widetilde{ss_2}$
uniformly random and unknown to $\mathcal{A}$ (it never appears in any
oracle output: the ciphertext $c_2^*$ is not in $\mathcal{O}_\Decap$'s
allowed inputs by IND-CCA2 rules), the only way $\mathcal{A}$
distinguishes $H(X)$ from uniform is to query $H$ at $X$. The probability of
a $q_H$-bounded adversary querying $H$ at $X$ when $X$ contains
$\widetilde{ss_2}$ of $\ell_2 \ge 256$ bits of min-entropy is at most
$q_H / 2^{\ell_2} \le q_H / 2^n$. Hence
\[
|\Pr[\Game_2 \to 1] - \Pr[\Game_3 \to 1]| \le \frac{q_H}{2^n}.
\]

\paragraph{In $\Game_3$, $ss^*$ is uniform.}
Therefore $\Pr[\Game_3 \to 1] = \frac{1}{2}$.

\paragraph{Combining.}
\begin{align*}
\Adv^{\mathsf{ind\text{-}cca}}_{\zwing}(\mathcal{A}) & = 2|\Pr[\Game_0 \to 1] - 1/2| \\
& \le 2 \sum_{j=0}^{2} |\Pr[\Game_j \to 1] - \Pr[\Game_{j+1} \to 1]| \\
& \le 2\!\left(\frac{q_H q_D}{2^n} + \Adv^{\mathsf{ind\text{-}cca}}_{\KEM_2}(\mathcal{B}_2) + \frac{q_H}{2^n}\right) \\
& \le 2 \Adv^{\mathsf{ind\text{-}cca}}_{\KEM_2}(\mathcal{B}_2) + \frac{2(q_H + q_H q_D)}{2^n}.
\end{align*}
By symmetry, the analogous bound holds with $i^\star = 1$ and $i^\star = 3$.
Taking the minimum and absorbing constants gives the theorem statement
(the additive correctness term comes from the small probability that any
of the three component KEMs decapsulates incorrectly during the simulation).
\qedhere
\end{proof}

\begin{remark}[On QROM]
The proof above uses classical random-oracle reasoning. For the post-quantum
threat model, we lift the analysis to the QROM via standard techniques: each
random-oracle game hop carries an additional $\mathcal{O}(q_H^2 / 2^n)$
multiplicative factor from the one-way-to-hiding (O2H) lemma. The resulting
bound retains the same shape but with worse constants. In particular, the
underlying $\KEM_i$ advantages must be the QROM-IND-CCA2 advantages, which
is exactly the form in which $\mlkem$ has been analysed.
\end{remark}

\section{Theorem 2: ZAP IND-CCA2-AKE with Z-Wing}\label{sec:thm2}

\subsection{The handshake}

For completeness we restate the relevant fragment of the \zap{} handshake.
Let $A, B$ be peers with long-term \zwing{} keys $(pk_A, sk_A), (pk_B, sk_B)$.
A session proceeds as:
\begin{enumerate}
\item $A \to B$: $(pk_A, c^A_e \gets \zwing.\Encap(pk_B), n_A)$, where
$n_A \xleftarrow{\$} \BB^{256}$ is a session nonce.
\item $B \to A$: $(pk_B, c^B_e \gets \zwing.\Encap(pk_A), n_B)$.
\item Both compute $ss_A \gets \zwing.\Decap(sk_A, c^B_e)$ and
$ss_B \gets \zwing.\Decap(sk_B, c^A_e)$, and likewise for the ephemeral
\zwing{} contributions.
\item Transport key:
$K_{\mathsf{tx}} \gets \hkdf(ss_A \parallel ss_B,\ \tau)$,
where $\tau = \texttt{"ZAP-v1"} \parallel \mathsf{fp}(pk_A) \parallel \mathsf{fp}(pk_B) \parallel c^A_e \parallel c^B_e \parallel n_A \parallel n_B \parallel \mathsf{suite}$.
\end{enumerate}

\subsection{The theorem}

\begin{theorem}[\zap{} with \zwing{} IND-CCA2-AKE]\label{thm:ake}
For any QPT adversary $\mathcal{A}$ in the AKE game with $n_P$ parties,
$n_S$ sessions per party, $q_H$ random-oracle queries, and $q_D$ decap
queries,
\[
\Adv^{\mathsf{ake}}_{\zap[\zwing]}(\mathcal{A}) \le
2 n_P n_S \cdot \Adv^{\mathsf{ind\text{-}cca}}_{\zwing}(\mathcal{B}) +
\frac{n_P^2 n_S^2}{2^{|\mathsf{nonce}|}} +
\frac{q_H^2}{2^{n+1}},
\]
where $\mathcal{B}$ runs in time $\approx t(\mathcal{A})$ plus simulation
overhead.
\end{theorem}

\begin{proof}
We use a hop sequence $\Game_0, \Game_1, \Game_2, \Game_3, \Game_4$.

\paragraph{$\Game_0$: Real AKE game.}
$\Pr[\Game_0 \to 1] = \frac{1}{2} + \frac{1}{2} \Adv^{\mathsf{ake}}(\mathcal{A})$.

\paragraph{$\Game_1$: Nonce-uniqueness abort.}
The simulator aborts if any two honest sessions sample the same nonce
$n \in \BB^{256}$.

\emph{Bound:} Birthday-bound over $n_P n_S$ sessions:
\[
|\Pr[\Game_0 \to 1] - \Pr[\Game_1 \to 1]| \le \binom{n_P n_S}{2} 2^{-256} \le \frac{n_P^2 n_S^2}{2^{|\mathsf{nonce}|+1}}.
\]

\paragraph{$\Game_2$: Random-oracle collision abort.}
Abort if any two distinct $H$-inputs collide in their image.

\emph{Bound:} $\binom{q_H}{2} 2^{-n} \le q_H^2 / 2^{n+1}$.

\paragraph{$\Game_3$: Test session guess.}
The simulator guesses, at the start of the experiment, the index $(P^\star, s^\star)$
of the session $\mathcal{A}$ will Test. If the guess is wrong, the simulator
aborts with output a coin-flip.

\emph{Bound:} The probability of guessing correctly is $1/(n_P n_S)$, so
\[
\Pr[\Game_3 \to 1] - \frac{1}{2} = \frac{1}{n_P n_S}\!\left(\Pr[\Game_2 \to 1] - \frac{1}{2}\right).
\]

\paragraph{$\Game_4$: Replace $\zwing{}$ shared secret in test session by
uniform.}
For the (guessed) test session $(P^\star, s^\star)$, the simulator
replaces the \zwing{}-derived $ss$ contributing to $K_{\mathsf{tx}}$ by a
uniformly random value of the same length.

\emph{Bound by Theorem~\ref{thm:zw-cca}:} The freshness condition guarantees
that for some $i^\star \in \{1, 2, 3\}$, the corresponding component KEM
secret in the test session is not corrupted. We use that
particular reduction (one of the three symmetric branches) to construct
$\mathcal{B}$:
\begin{itemize}
\item $\mathcal{B}$ receives a $\KEM_{i^\star}$ challenge $(pk_{i^\star}, c^*_{i^\star}, ss^*_{i^\star})$.
\item $\mathcal{B}$ embeds $pk_{i^\star}$ as the corresponding component of the
target party's long-term \zwing{} public key (or ephemeral, as appropriate),
and uses $c^*_{i^\star}$ as the matching component of the test-session
ciphertext.
\item For all other components and other sessions, $\mathcal{B}$ generates
keys honestly and answers all corruption / decap / session-key queries by
direct computation.
\item For decap queries against the embedded $pk_{i^\star}$, $\mathcal{B}$
uses its own $\KEM_{i^\star}$ decapsulation oracle.
\end{itemize}
The simulation perfectly matches $\Game_3$ if $ss^*_{i^\star}$ is real, and
matches $\Game_4$ if random. By the IND-CCA2 reduction,
\[
|\Pr[\Game_3 \to 1] - \Pr[\Game_4 \to 1]| \le \Adv^{\mathsf{ind\text{-}cca}}_{\zwing}(\mathcal{B}).
\]

\paragraph{In $\Game_4$, $K_{\mathsf{tx}}$ is uniform.}
By Lemma~\ref{lem:transcript}, $K_{\mathsf{tx}} = \hkdf(ss, \tau)$ is uniform
in $\BB^{|K|}$ when $ss$ is uniform and $\hkdf$ is modelled as a random
oracle, except for the collision probability already accounted for in
$\Game_2$.

So $\Pr[\Game_4 \to 1] = 1/2$.

\paragraph{Combining.}
\begin{align*}
\Adv^{\mathsf{ake}}(\mathcal{A}) & = 2 \left(\Pr[\Game_0 \to 1] - 1/2\right) \\
& \le 2 \cdot n_P n_S \left(\Pr[\Game_3 \to 1] - 1/2\right) + \frac{n_P^2 n_S^2}{2^{|\mathsf{nonce}|}} + \frac{q_H^2}{2^n} \\
& \le 2 n_P n_S \cdot \Adv^{\mathsf{ind\text{-}cca}}_{\zwing}(\mathcal{B}) + \frac{n_P^2 n_S^2}{2^{|\mathsf{nonce}|}} + \frac{q_H^2}{2^n},
\end{align*}
which proves the theorem. \qedhere
\end{proof}

\section{Hop-Sequence Summary}\label{sec:summary}

\begin{table}[h]
\centering
\caption{\zwing{} IND-CCA2 hop sequence (case $i^\star = 2$).}
\begin{tabular}{lll}
\hline
Game & Difference from previous & Bound on transition \\
\hline
$\Game_0$ & Real IND-CCA2 game & --- \\
$\Game_1$ & Decap-oracle programming abort & $q_H q_D / 2^n$ \\
$\Game_2$ & $ss_2^*$ replaced by uniform & $\Adv^{\mathsf{ind\text{-}cca}}_{\KEM_2}(\mathcal{B}_2)$ \\
$\Game_3$ & $ss^*$ replaced by uniform & $q_H / 2^n$ \\
\hline
\end{tabular}
\end{table}

\begin{table}[h]
\centering
\caption{\zap{} with \zwing{} IND-CCA2-AKE hop sequence.}
\begin{tabular}{lll}
\hline
Game & Difference & Bound \\
\hline
$\Game_0$ & Real AKE game & --- \\
$\Game_1$ & Nonce uniqueness & $n_P^2 n_S^2 / 2^{|\mathsf{nonce}|+1}$ \\
$\Game_2$ & RO collision avoidance & $q_H^2 / 2^{n+1}$ \\
$\Game_3$ & Test-session guess & factor $n_P n_S$ in advantage \\
$\Game_4$ & \zwing{} ss in test by uniform & $\Adv^{\mathsf{ind\text{-}cca}}_{\zwing}(\mathcal{B})$ \\
\hline
\end{tabular}
\end{table}

\section{Concrete Bound}\label{sec:concrete}

For deployment parameters $n = 256$, $|\mathsf{nonce}| = 256$,
$n_P \le 2^{20}$, $n_S \le 2^{20}$, $q_H \le 2^{60}$, $q_D \le 2^{40}$,
\[
\Adv^{\mathsf{ake}}_{\zap[\zwing]}(\mathcal{A}) \le
2^{41} \cdot \Adv^{\mathsf{ind\text{-}cca}}_{\zwing}(\mathcal{B}) +
2^{40} \cdot 2^{-256} + 2^{120} \cdot 2^{-257}
\le 2^{41} \cdot \min_i \Adv^{\mathsf{ind\text{-}cca}}_{\KEM_i} + 2^{-136}.
\]
With \mlkem{}-768 conjectured at $\approx 2^{-128}$ classical advantage and
\xtf{} at $\approx 2^{-128}$ for the elliptic-curve component, the AKE advantage
is bounded by $2^{-87}$ in the multi-user multi-session setting, which is
comfortably below any practical threshold. The dominant term is the
guess-the-test-session factor; tighter multi-user reductions for the
underlying KEMs would directly translate.

\section{Forward Secrecy and Post-Compromise Security}\label{sec:fs}

The hybrid freshness condition (\cref{def:fresh}) yields
\emph{weak forward secrecy}: an adversary that corrupts both peers' static
keys \emph{after} a session completes still cannot recover the session key
provided neither peer's ephemeral was revealed and at least one ephemeral
component remains uncorrupted. This is because the test session's ephemeral
$\KEM_1$ keypair is short-lived and securely erased.

\emph{Strong} forward secrecy --- security against post-session corruption
of \emph{ephemerals} also --- is unattainable in the absence of additional
key-rotation machinery, and is not claimed.

\emph{Post-compromise security} after key rotation: \zap{} performs handshake
re-keying every $T$ messages or $\Delta t$ seconds. A new \zwing{} encap on
a fresh ephemeral disconnects the post-rotation key from the pre-rotation
KEM secrets, except via $\sha$ collisions in the transcript, bounded by
Lemma~\ref{lem:transcript}.

\section{Open Questions Recap}\label{sec:open}

The companion paper lists open questions; we add here:
\begin{itemize}
\item \emph{Tightness.} The factor $n_P n_S$ in Theorem~\ref{thm:ake} is from
guessing the test session. A tight reduction would use the multi-user
techniques of Bader--Jager--Li--Sch{\"a}ge.
\item \emph{QROM bounds.} Our hop sequence is classical-ROM. A QROM-tight
analysis of \zwing{} via the O2H framework is the obvious follow-up.
\item \emph{Mechanisation.} The proof structure is suitable for an
EasyCrypt mechanisation; the only delicate step is the simultaneous
embedding of three KEM challenges, which requires care with the simulator's
random-oracle programming.
\end{itemize}

\appendix

\section{Bibliographic Notes}\label{app:bib}

The sequence-of-games style here follows Shoup \cite{shoup2004games}.
The robust-combiner construction generalises the parallel-encap-then-hash
construction of \cite{giacon2018kem}, with the public-key binding of
\xwing{} \cite{xwing2024}. The hybrid freshness condition is from
\cite{bindel2019hybrid}. The HKDF random-oracle abstraction is conservative;
the original analysis \cite{krawczyk2010hkdf} relies on a weaker
``computational extractor'' assumption.

\begin{thebibliography}{9}
\bibitem{xwing2024} Barbosa, Connolly, Duarte, Kaiser, Schwabe, Varner, Westerbaan, H{\"u}lsing.
``X-Wing: The Hybrid KEM You've Been Looking For.'' eprint 2024/039.
\bibitem{kyber2018} Bos, Ducas, Kiltz, Lepoint, Lyubashevsky, Schanck, Schwabe, Seiler, Stehl{\'e}.
``CRYSTALS-Kyber: A CCA-Secure Module-Lattice-Based KEM.'' EuroS\&P 2018.
\bibitem{nistmlkem} NIST. ``FIPS 203: Module-Lattice-Based Key-Encapsulation Mechanism Standard.'' 2024.
\bibitem{giacon2018kem} Giacon, Heuer, Poettering. ``KEM Combiners.'' PKC 2018.
\bibitem{bindel2019hybrid} Bindel, Brendel, Fischlin, Goncalves, Stebila.
``Hybrid Key Encapsulation Mechanisms and Authenticated Key Exchange.'' PQCrypto 2019.
\bibitem{krawczyk2010hkdf} Krawczyk. ``Cryptographic Extraction and Key Derivation: The HKDF Scheme.'' CRYPTO 2010.
\bibitem{shoup2004games} Shoup. ``Sequences of Games: A Tool for Taming Complexity in Security Proofs.'' eprint 2004/332.
\bibitem{cremers2017ake} Cremers, Horvat, Hoyland, Scott, van der Merwe. ``A Comprehensive Symbolic Analysis of TLS 1.3.'' CCS 2017.
\bibitem{boyd2020ake} Boyd, Mathuria, Stebila. \emph{Protocols for Authentication and Key Establishment}, 2nd ed., Springer 2020.
\end{thebibliography}

\end{document}
